\section{Specializing the original finite free module complex}

\begin{lemma}[The original tensor extension retains its matrix]
\label{lem:kn-specialization-original-base-change-matrix}
\lean{EnvelopeProperty.Vanishing.SpecializationOriginalFreeModules.toMatrix_baseChange}
\leanok
Let $f:A\to B$ be a linear map between finite free $R$-modules with
chosen finite bases.  Over an $R$-algebra $S$, its actual tensor
extension $1\otimes f$ has, in the base-changed bases, the matrix
obtained by applying the original coefficient map $R\to S$ to every
entry of the original matrix of $f$.
\end{lemma}
\begin{proof}
The $j$th base-changed basis vector is $1\otimes a_j$.  Its image
is $1\otimes f(a_j)$.  The coefficient in the $i$th base-changed
basis vector is the image in $S$ of the original $i$th coefficient
of $f(a_j)$.  This proves equality of every matrix entry.
\end{proof}

\begin{lemma}[Exactness retains the original basis identifications]
\label{lem:kn-specialization-original-exactness-matrix}
\lean{EnvelopeProperty.Vanishing.SpecializationOriginalFreeModules.matrix_exact_iff}
\leanok
For an original pair of linear maps between finite free modules,
exactness of its two matrices over the same coefficient ring is equivalent to exactness
of the original pair of linear maps under the original basis
identifications.
\end{lemma}
\begin{proof}
The matrix maps are obtained by conjugating the original maps with
the coordinate isomorphisms.  Precomposition by an isomorphism
preserves the image; postcomposition by an isomorphism preserves
the zero locus.  The common middle coordinate isomorphism transports
the range of the first map and the kernel of the second to the
corresponding original range and kernel.  Their equality is therefore
equivalent in the two representations.
\end{proof}

\begin{theorem}[A principal exactness neighborhood for the original free modules]
\label{thm:kn-specialization-original-free-neighborhood}
\lean{EnvelopeProperty.Vanishing.SpecializationOriginalFreeModules.exists_closed_fibre_exactness_neighborhood}
\leanok
\uses{lem:kn-specialization-original-base-change-matrix,lem:kn-specialization-original-exactness-matrix,thm:kn-specialization-finite-free-exactness}
Let $A\xrightarrow f B\xrightarrow g C$ be an original finite free
$R$-module complex and let $\mathfrak p$ be an original maximal
ideal.  If the actual tensor maps over $\kappa(\mathfrak p)$ are
exact, there is $r\notin\mathfrak p$ such that the actual tensor
maps over every original $R$-algebra field in which $r$ is nonzero
are exact.
\end{theorem}
\begin{proof}
Choose finite bases of the original free modules.  The first two
lemmas identify exactness of the actual residue tensor maps with
exactness of the original matrices after the original residue
coefficient map.  The matrix specialization theorem supplies an
element $r$ with residue $1$ and exactness at every field point
where it remains nonzero.  The same two lemmas identify these matrix
maps with the actual tensor maps of the original modules.  Thus no
matrix-comparison isomorphism is a supplied hypothesis.
\end{proof}

\begin{corollary}[Generic exactness for the original tensor complex]
\label{cor:kn-specialization-original-free-generic}
\lean{EnvelopeProperty.Vanishing.SpecializationOriginalFreeModules.generic_exact_of_closed_fibre_exact}
\leanok
\uses{thm:kn-specialization-original-free-neighborhood}
If $R$ is a domain, exactness of the actual closed residue tensor
complex implies exactness of the actual tensor extension of the
same original finite free module complex to any fraction field of $R$.
\end{corollary}
\begin{proof}
The element $r\notin\mathfrak p$ supplied by the theorem is
nonzero in $R$.  Injectivity of the canonical map to the fraction
field retains nonvanishing of $r$, so the theorem applies to that
original canonical coefficient map.
\end{proof}

\begin{corollary}[Native homology of the original tensor complex]
\label{cor:kn-specialization-original-free-native-homology}
\lean{EnvelopeProperty.Vanishing.SpecializationOriginalFreeModules.exists_original_homology_zero_neighborhood,EnvelopeProperty.Vanishing.SpecializationOriginalFreeModules.original_generic_homology_isZero_of_closed_fibre}
\leanok
\uses{thm:kn-specialization-original-free-neighborhood,cor:kn-specialization-original-free-generic}
Form the native short complex from the actual tensor extensions of the
same original maps.  If its native homology vanishes over the canonical
closed residue fibre, it vanishes after every coefficient map to a
field in one constructed principal neighborhood.  Over a domain it
also vanishes after the canonical fraction-field extension.
\end{corollary}
\begin{proof}
The native homology of a module short complex vanishes exactly when
the original range equals the original kernel.  Apply this criterion
to the actual residue tensor maps, the principal-neighborhood
exactness theorem, and the same criterion in the reverse direction.
For the generic field use its canonical coefficient map and the
generic exactness corollary.
\end{proof}
